%!TeX root = ./../Presentation.tex

\section{Problemstellung und Zielstellung}

\begin{frame}[fragile]
    \begin{center}
        \huge\textbf{\insertsection}
    \end{center}
\end{frame}

\begin{frame}[fragile]
\frametitle{\insertsection}
% Was ist das Problem?

\begin{itemize}
    \item Inkompatibilität der Befehlssatzarchitekturen zwischen Entwicklungs-
        und Zielplattform
    \item Lieferengpässe bei Halbleitertechnologie aufgrund externer Faktoren
    \item Mangelhafter Zugang zu Testhardware
\end{itemize}
\vspace{5mm}
\hspace{5mm}\rightarrow\hspace{2mm} Reduktion der Entwicklungsgeschwindigkeit
\end{frame}

\begin{frame}[fragile]
\frametitle{\insertsection}
% Wie löse ich Problem

\begin{itemize}
    \item Emulation des Mikrocontrollers \textbf{STM32F429} mittels QEMU
    \item Ausführung der Software für eingebettete Systeme
    \begin{itemize}
        \item in emulierter Umgebung
        \item auf realer Hardware-Umgebung
    \end{itemize}
\end{itemize}
\begin{alertblock}{Problem}
    QEMU enthält keine Integration für den STM32F429 Mikrocontroller,
\end{alertblock}
\vspace{3mm}
aber Integrationen für verwandete Mikrocontroller der STM32F2 und F4
Familien und kompatibler Peripherie!

\end{frame}

\begin{frame}[fragile]
\frametitle{\insertsection}
\framesubtitle{Ablauf der Implementierung}

\begin{enumerate}
    \item Implementierung der Maschine und des SoC
    \item Identifikation geeigneter, kompatibler Peripherie
    \item prototypische Implementierung der Peripherie in QEMU
    \begin{itemize}
        \item direkt in QEMU als Teil des Mikrocontrollers
        \item als externer Prozess mittels EDI
    \end{itemize}
    \item Test der Implementierungen in QEMU und auf Hardware
\end{enumerate}

\end{frame}
